\documentclass[11pt,a4paper]{article}

\usepackage[T1]{fontenc}
\usepackage[utf8]{inputenc}
\usepackage[french]{babel}
\usepackage{amsmath,amssymb,amsthm}
\usepackage{tikz}
\usepackage{tikz-cd}
\usetikzlibrary{arrows.meta,positioning,decorations.pathreplacing,patterns,calc,shapes.geometric}
\usepackage{booktabs}
\usepackage{geometry}
\geometry{margin=2.4cm}
\usepackage{hyperref}
\hypersetup{colorlinks=true,linkcolor=black,urlcolor=blue}

\theoremstyle{plain}
\newtheorem{theoreme}{Théorème}
\newtheorem{forme}[theoreme]{Forme}
\theoremstyle{definition}
\newtheorem{defi}[theoreme]{Définition}
\theoremstyle{remark}
\newtheorem{obs}[theoreme]{Observation}

\newcommand{\Darb}{\mathsf{Darb}}
\newcommand{\Cont}{\mathsf{Cont}}
\newcommand{\Der}{\mathsf{Der}}
\newcommand{\Inj}{\mathsf{Inj}}
\newcommand{\Mono}{\mathsf{Mono}}
\newcommand{\CountIm}{\mathsf{CountIm}}
\newcommand{\Const}{\mathsf{Const}}
\newcommand{\SMA}{\mathsf{StrictMonoAnti}}
\newcommand{\R}{\mathbb{R}}

\title{Darboux : raisonnement diagrammatique,\\ catégories, types, géométrie\\[2mm]
\large et une carte de l'espace de preuves (formalisée en Lean 4)}
\author{}
\date{}

\begin{document}
\maketitle

\begin{abstract}
\noindent Prose minimale. Diagrammes, tables, énoncés. Chaque énoncé marqué
\textsf{[Lean]} est démontré sans \texttt{sorry} dans \texttt{RequestProject/Darboux/}.
\end{abstract}

\section{Deux Darboux}

\begin{center}
\begin{tabular}{@{}lll@{}}
\toprule
& \textbf{Darboux (analyse)} & \textbf{Darboux (symplectique)}\\
\midrule
Objet & $f' : I \to \R$ & $\omega$ forme symplectique sur $M^{2n}$\\
Énoncé & $f'(I)$ est un intervalle & $\exists$ carte : $\omega = \sum dp_i\wedge dq_i$\\
Slogan & pas de saut & pas d'invariant local\\
Mécanisme & connexité $+$ Fermat & astuce de Moser (chemin $\omega_t$)\\
Obstruction levée & continuité de $f'$ & modules locaux\\
Dimension $1$ & --- & normalisation d'une densité \textsf{[Lean]}\\
\bottomrule
\end{tabular}
\end{center}

\medskip
\noindent Le pont : les deux énoncés disent \emph{\og{}l'image d'un connexe est connexe\fg{}} pour
un objet qui n'est pas une application continue (une dérivée ; un germe de forme).

\section{Visualisations}

\subsection{Ce qui est interdit : le saut}

\begin{center}
\begin{tikzpicture}[scale=1.05]
\begin{scope}
  \draw[->] (-0.3,0) -- (3.4,0) node[right]{$x$};
  \draw[->] (0,-0.3) -- (0,2.4) node[above]{$f'$};
  \draw[very thick] (0.2,0.4) -- (1.5,0.4);
  \draw[very thick] (1.5,1.8) -- (2.9,1.8);
  \fill (1.5,0.4) circle (1.4pt);
  \draw[dashed] (0,1.1) -- (3.1,1.1) node[right]{$m$};
  \node[red] at (1.5,1.1) {\Large $\times$};
  \node at (1.55,-0.55) {\small saut : $m$ jamais atteint};
  \node at (1.55,-1.0) {\small \textbf{impossible} pour une dérivée};
\end{scope}
\begin{scope}[xshift=6cm]
  \draw[->] (-0.3,0) -- (3.4,0) node[right]{$x$};
  \draw[->] (0,-0.3) -- (0,2.4) node[above]{$f'$};
  \draw[very thick,domain=0.2:2.9,samples=200,smooth]
    plot (\x,{1.1+0.75*sin(360*\x/0.55)*exp(-0.15*\x)});
  \draw[dashed] (0,1.1) -- (3.1,1.1) node[right]{$m$};
  \node at (1.55,-0.55) {\small oscillation sauvage : \textbf{permise}};
  \node at (1.55,-1.0) {\small ($f'$ non continue, mais $\Darb$)};
\end{scope}
\end{tikzpicture}
\end{center}

\noindent Témoin : $f(x)=x^{2}\sin(1/x)$, $f(0)=0$ ; $f'$ existe partout, n'est pas continue
en $0$, et vérifie la propriété des valeurs intermédiaires.

\subsection{La preuve classique en une image (Fermat)}

\begin{center}
\begin{tikzpicture}[scale=1.1]
  \draw[->] (-0.3,0) -- (5.2,0) node[right]{$x$};
  \draw[->] (0,-1.5) -- (0,1.6) node[above]{$g=f-m\,x$};
  \draw[very thick,domain=0.3:4.7,samples=200,smooth]
    plot (\x,{0.9*(\x-0.3)*(\x-4.7)*0.28});
  \fill (0.3,0) circle (1.6pt) node[above left]{$a$};
  \fill (4.7,0) circle (1.6pt) node[above right]{$b$};
  \draw[dotted] (2.5,-1.11) -- (2.5,0) node[above]{$c$};
  \fill (2.5,-1.11) circle (1.6pt);
  \draw[<-] (2.9,-1.25) -- (4.0,-1.25) node[right]{\small $g'(c)=0$};
  \draw[-{Stealth}] (0.35,0.05) -- (0.9,-0.25) node[midway,above right]{\small $g'(a)<0$};
  \draw[-{Stealth}] (4.65,0.05) -- (4.1,-0.25) node[midway,above left]{\small $g'(b)>0$};
\end{tikzpicture}
\end{center}

\noindent $f'(a)<m<f'(b)$ $\Rightarrow$ $g=f-mx$ a un minimum \emph{intérieur} $c$
$\Rightarrow$ $g'(c)=0$ $\Rightarrow$ $f'(c)=m$. Ingrédients : compacité et Fermat.

\subsection{Configuration à trois points (injectivité)}

\begin{center}
\begin{tikzpicture}[scale=1.0]
  \draw[->] (-0.3,0) -- (5.6,0) node[right]{$x$};
  \draw[->] (0,-0.3) -- (0,3.1) node[above]{$f$};
  \fill (0.6,0.7) circle (1.6pt) node[below right]{$(a,f(a))$};
  \fill (2.6,2.6) circle (1.6pt) node[above]{$(b,f(b))$};
  \fill (4.8,0.9) circle (1.6pt) node[below left]{$(c,f(c))$};
  \draw[dashed] (0,1.9) -- (5.3,1.9) node[right]{$z$};
  \draw[thick,gray] plot[smooth,tension=0.8] coordinates {(0.6,0.7) (1.6,2.0) (2.6,2.6)};
  \draw[thick,gray] plot[smooth,tension=0.8] coordinates {(2.6,2.6) (3.7,1.7) (4.8,0.9)};
  \fill[red] (1.52,1.9) circle (1.8pt) node[below]{\small $p$};
  \fill[red] (3.52,1.9) circle (1.8pt) node[below]{\small $q$};
  \node at (2.6,-0.75) {\small $f(p)=f(q)=z$, $p\neq q$ : contredit l'injectivité};
\end{tikzpicture}
\end{center}

\noindent Si $f(b)$ est un extremum strict du triplet, la valeur intermédiaire $z$ est
atteinte deux fois : c'est le c\oe ur de la Forme~\ref{f:inj}.

\subsection{Deux autres formes, en image}

\begin{center}
\begin{tikzpicture}[scale=0.95]
\begin{scope}
  \draw[->] (-0.2,0) -- (4.4,0) node[right]{$x$};
  \draw[->] (0,-0.2) -- (0,2.8) node[above]{$f$};
  \draw[very thick] (0.3,0.5) -- (1.9,0.5);
  \draw[very thick] (1.9,2.0) -- (4.0,2.0);
  \draw[decorate,decoration={brace,amplitude=5pt},thick]
    (4.25,0.5) -- (4.25,2.0) node[midway,right]{\small trou};
  \node at (2.1,-0.7) {\small $\Mono$ : un saut $\Rightarrow$ un trou dans l'image};
  \node at (2.1,-1.15) {\small donc $\Darb\ \&\ \Mono \Rightarrow \Cont$};
\end{scope}
\begin{scope}[xshift=7.2cm]
  \draw[->] (-0.2,0) -- (4.4,0) node[right]{$x$};
  \draw[->] (0,-0.2) -- (0,2.8) node[above]{$f$};
  \foreach \y in {0.4,1.0,1.6,2.2} \draw[dotted] (0.2,\y) -- (4.1,\y);
  \draw[very thick] (0.3,1.0) -- (4.0,1.0);
  \node at (2.1,-0.7) {\small image dénombrable $\subseteq\{y_0,y_1,\dots\}$};
  \node at (2.1,-1.15) {\small un intervalle non trivial est indénombrable};
\end{scope}
\end{tikzpicture}
\end{center}

\subsection{La carte : l'espace de preuves}

\begin{center}
\begin{tikzpicture}[
  nd/.style={draw,rounded corners=2pt,inner sep=3pt,font=\small},
  hy/.style={circle,fill=black,inner sep=1.3pt},
  ar/.style={-{Stealth[length=2mm]}},
  hyar/.style={-{Stealth[length=2mm]},densely dashed}]

  \node[nd] (der)  at (0,3)     {$\Der$};
  \node[nd] (cont) at (0,1.5)   {$\Cont$};
  \node[nd] (darb) at (3,2.25)  {$\Darb$};
  \node[nd] (inj)  at (3,0.2)   {$\Inj$};
  \node[nd] (mono) at (0,0)     {$\Mono$};
  \node[nd] (cnt)  at (6.8,3.3) {$\CountIm$};
  \node[nd] (cst)  at (9.6,2.2) {$\Const$};
  \node[nd] (sma)  at (6.8,0.2) {$\SMA$};
  \node[nd] (cg)   at (3,-1.5)  {$\mathsf{ClosedGraph}$};

  \draw[ar] (der) -- (darb);
  \draw[ar] (cont) -- (darb);
  \draw[ar] (cst) to[bend right=20] (cont);
  \draw[ar] (cst) to[bend left=32] (mono);
  \draw[ar] (cst) to[bend right=12] (cnt);
  \draw[ar] (sma) -- (inj);

  \node[hy] (h1) at (5.0,1.1) {};
  \draw (darb) -- (h1); \draw (inj) -- (h1);
  \draw[hyar] (h1) -- (sma);

  \node[hy] (h2) at (5.4,2.9) {};
  \draw (darb) -- (h2); \draw (cnt) -- (h2);
  \draw[hyar] (h2) -- (cst);

  \node[hy] (h3) at (1.6,1.0) {};
  \draw (darb) -- (h3); \draw (mono) -- (h3);
  \draw[hyar] (h3) to[bend left=10] (cont);

  \node[hy] (h4) at (1.2,-0.7) {};
  \draw (darb) -- (h4); \draw (cg) -- (h4);
  \draw[hyar] (h4) to[bend right=12] (cont);
  \draw[ar] (cont) to[bend right=18] (cg);

  \node[font=\scriptsize,anchor=west] at (-1.2,-2.4)
    {\textbf{---} arête prouvée \qquad \textbf{-\,-} hyper-arête (deux hypothèses) \qquad
     $\bullet$ n\oe ud de conjonction};
\end{tikzpicture}
\end{center}

\noindent $\Const$ est le puits, $\Darb$ le carrefour. Non-arêtes (obstructions vraies) :
$\Darb \nRightarrow \Cont$ ; $\Darb$ n'est stable ni par somme ni par composition ;
$\Der \nRightarrow \Cont$.

\subsection{Post-composition : la seule stabilité}

\begin{center}
\begin{tikzcd}[column sep=large,row sep=large]
s \arrow[r,"f\ \in\ \Darb"] \arrow[dr,"g\circ f\ \in\ \Darb"',dashed]
  & f(s) \arrow[d,"g\ \in\ \Cont"]\\
  & g(f(s))
\end{tikzcd}
\end{center}

\noindent \textsf{[Lean]} \texttt{DarbouxOn.comp\_continuousOn}. Rôles échangés
($f\in\Cont$, $g\in\Darb$), la conclusion tombe. $\Darb$ n'est donc pas une catégorie : c'est
un module à gauche sur $\Cont$.

\section{Catégories}

\begin{center}
\begin{tabular}{@{}ll@{}}
\toprule
Structure & Lecture de Darboux\\
\midrule
$\pi_0 : \mathbf{Top}\to\mathbf{Set}$ & $\Darb$ = \og{}préserve $\pi_0$\fg{} sans être un morphisme\\
$\mathbf{OC}$ & objets : parties ordre-connexes de $\R$\\
$\Darb(s)$ & module sur $\Cont$ : $g\cdot f := g\circ f$\\
Multicatégorie & hyper-arêtes = multiflèches $(\Darb,\Inj)\to \SMA$\\
Faisceaux & Darboux symplectique = trivialité locale du faisceau des $\omega$\\
Pseudogroupe & $\mathrm{Symp}$ agit transitivement sur les germes\\
\bottomrule
\end{tabular}
\end{center}

\begin{center}
\begin{tikzcd}[column sep=huge,row sep=large]
\{\text{parties ordre-connexes}\} \arrow[r,"f_*"] \arrow[d,"\pi_0"']
  & \{\text{parties ordre-connexes}\} \arrow[d,"\pi_0"]\\
\{*\} \arrow[r,equal] & \{*\}
\end{tikzcd}
\end{center}

\noindent Lecture : $f$ de Darboux $\iff$ $f_*$ préserve la fibre \og{}connexe\fg{} de $\pi_0$.
Continuité $\Rightarrow$ fonctorialité ; Darboux $\Rightarrow$ préservation seule, sans
fonctorialité (pas de stabilité par composition).

\smallskip
\noindent Version opéradique : un lemme à deux hypothèses est une opération
$\mu : \Darb\otimes X \to Y$. La carte de la \S2.5 est le graphe sous-jacent d'une
multicatégorie libre engendrée par les lemmes prouvés.

\section{Types}

\begin{center}
\begin{tabular}{@{}ll@{}}
\toprule
Énoncé & Type\\
\midrule
$\Darb$ & $\Pi_{t\subseteq s}\ \mathrm{OC}(t)\to \mathrm{OC}(f(t))$\\
valeur intermédiaire & $\Pi_{a,b\in s}\Pi_{y\in[f a,f b]}\ \Sigma_{c\in[a,b]}\ f(c)=y$\\
hyper-arête & $\Darb \to \Inj \to \SMA$ (règle à deux prémisses)\\
non-arête & type \emph{non} habité (contre-exemple)\\
\bottomrule
\end{tabular}
\end{center}

\begin{center}
$\dfrac{\ \Gamma\vdash h_1:\Darb(f,s)\qquad \Gamma\vdash h_2:\Inj(f,s)\ }
{\ \Gamma\vdash \mathsf{sma}\,h_1\,h_2:\SMA(f,s)\ }\quad\textsc{(forme~\ref{f:inj})}$
\end{center}

\noindent Contenu constructif : le $\Sigma$ ci-dessus n'est pas réalisable
constructivement (le théorème des valeurs intermédiaires classique implique LLPO) ; la
preuve formalisée passe par le choix classique. Version constructive usuelle :
$\forall\varepsilon>0,\ \exists c,\ |f(c)-y|<\varepsilon$.

\smallskip
\noindent Lecture \og{}recherche de preuves\fg{} : un lemme non prouvé est un \emph{trou}, un type
dont on cherche un habitant. \texttt{proofspace.py} énumère les trous ;
\texttt{Space.lean} certifie ceux qui sont comblés.

\section{Géométrie}

\subsection{Moser en une image}

\begin{center}
\begin{tikzpicture}[scale=1]
  \draw[rounded corners=4pt,thick] (-5.4,-2.0) rectangle (5.4,2.1);
  \node[anchor=north] at (0,2.05) {\small espace des germes de formes symplectiques en $p$};
  \node (w0) at (-3.6,0.1) {$\omega_0=\omega$};
  \node (w1) at (3.5,0.1) {$\omega_1=\sum dp_i\wedge dq_i$};
  \draw[very thick,-{Stealth}] (w0) to[bend left=22]
    node[midway,above]{\small $\omega_t=(1-t)\omega_0+t\omega_1$} (w1);
  \draw[dashed,-{Stealth}] (w0) to[bend right=22]
    node[midway,below]{\small flot $\psi_t$ de $X_t$, $\iota_{X_t}\omega_t=-\alpha$} (w1);
  \node at (0,-1.65)
    {$\psi_t^{*}\omega_t=\omega_0$ \quad$\Longrightarrow$\quad $\psi_1^{*}\omega_1=\omega_0$};
\end{tikzpicture}
\end{center}

\noindent Ingrédients : $\omega_1-\omega_0=d\alpha$ (lemme de Poincaré, local) ;
non-dégénérescence sur un voisinage ; l'équation $\iota_{X_t}\omega_t=-\alpha$ est
\emph{algébrique} en $X_t$. La rigidité vient de là.

\subsection{Dimension 1 : la densité, formalisée}

\begin{center}
\begin{tikzpicture}[scale=1]
\begin{scope}
  \draw[->] (-0.2,0) -- (3.6,0) node[right]{$x$};
  \draw[->] (0,-0.2) -- (0,2.2) node[above]{$\rho$};
  \draw[very thick,domain=0.15:3.3,samples=120,smooth]
    plot (\x,{0.6+0.5*\x-0.09*\x*\x});
  \node at (1.7,-0.7) {\small densité $\rho>0$};
\end{scope}
\draw[-{Stealth},very thick] (4.2,1.0) -- (5.6,1.0)
  node[midway,above]{\small $\Phi(x)=\int_0^x\rho$};
\begin{scope}[xshift=6.4cm]
  \draw[->] (-0.2,0) -- (3.6,0) node[right]{$u=\Phi(x)$};
  \draw[->] (0,-0.2) -- (0,2.2) node[above]{$1$};
  \draw[very thick] (0.15,1.0) -- (3.3,1.0);
  \node at (1.7,-0.7) {\small densité constante $du$};
\end{scope}
\end{tikzpicture}
\end{center}

\noindent \textsf{[Lean]} \texttt{Darboux.exists\_chart\_of\_density} :
$\Phi$ strictement croissante, $\Phi'=\rho$, $\mathrm{im}\,\Phi$ intervalle,
$\int_a^b\rho=\Phi(b)-\Phi(a)$. Moser en dimension~$1$ : la carte est la primitive, et son
image est un intervalle par le premier Darboux.

\section{Grandes idées}

\begin{enumerate}\itemsep2pt
\item \textbf{Connexité transportée sans continuité.} Darboux sépare \og{}préserver les
intervalles\fg{} de \og{}être continue\fg{}. Le n\oe ud $\Darb$, pas $\Der$, est l'unité d'analyse
qui fait marcher les preuves.
\item \textbf{Rigidité par cardinalité.} Un intervalle non trivial est indénombrable :
toute contrainte dénombrable sur les valeurs force la constance (Forme~\ref{f:count}).
\item \textbf{Rigidité par ordre.} Injectivité $+$ valeurs intermédiaires $=$ monotonie
stricte (Forme~\ref{f:inj}) ; monotonie $+$ valeurs intermédiaires $=$ continuité
(Forme~\ref{f:mono}).
\item \textbf{Ce que la linéarité ajoute.} $\Der$ est un espace vectoriel, $\Darb$ non :
d'où des formes réservées aux dérivées, comme $f'(x)\neq x$ partout $\Rightarrow$ signe
constant de $f'-\mathrm{id}$ (Forme~\ref{f:fix}).
\item \textbf{Pas d'invariant local} (versant symplectique) : \og{}pas de saut\fg{}, un cran plus
haut.
\item \textbf{Un espace de preuves est un hypergraphe.} Les conjectures non cartographiées
sont ses trous : paires ouvertes et triplets $(X,Y)\to Z$ non couverts.
\end{enumerate}

\section{Formes récoltées (toutes formalisées)}

Cadre commun : $s\subseteq\R$ ordre-connexe, $f:\R\to\R$.

\begin{forme}[valeur évitée]\label{f:sign}
$f\in\Darb(s)$ et $f\neq m$ sur $s$ $\Rightarrow$ $f<m$ partout, ou $f>m$ partout.
\par\smallskip\hfill\small\texttt{DarbouxOn.forall\_lt\_or\_forall\_gt}
\end{forme}

\begin{forme}[image dénombrable]\label{f:count}
$f\in\Darb(s)$ et $f(s)$ dénombrable $\Rightarrow$ $f$ constante sur $s$ ; en particulier
une dérivée à valeurs entières est constante.
\par\smallskip\hfill\small\texttt{DarbouxOn.subsingleton\_image\_of\_countable}, \texttt{deriv\_eq\_of\_int\_valued}
\end{forme}

\begin{forme}[injectivité]\label{f:inj}
$f\in\Darb(s)$ et $f$ injective sur $s$ $\Rightarrow$ $f$ strictement monotone sur $s$ ;
généralisation sans continuité de \texttt{ContinuousOn.strictMonoOn\_of\_injOn\_Icc'}.
\par\smallskip\hfill\small\texttt{DarbouxOn.strictMonoOn\_or\_strictAntiOn}
\end{forme}

\begin{forme}[monotonie]\label{f:mono}
$f\in\Darb(s)$, $f$ croissante, $s$ ouvert $\Rightarrow$ $f$ continue sur $s$.
\par\smallskip\hfill\small\texttt{DarbouxOn.continuousOn\_of\_monotoneOn\_of\_isOpen}
\end{forme}

\begin{forme}[convexe $\Rightarrow$ $C^1$]\label{f:cvx}
$f$ convexe et dérivable sur un ouvert convexe $\Rightarrow$ $f'$ continue (corollaire de
\ref{f:mono} : $f'$ est monotone \emph{et} de Darboux).
\par\smallskip\hfill\small\texttt{Darboux.ConvexOn.continuousOn\_deriv}
\end{forme}

\begin{forme}[point fixe des dérivées]\label{f:fix}
$f$ dérivable sur $s$ avec $f'(x)\neq x$ partout $\Rightarrow$ $f'<\mathrm{id}$ partout ou
$f'>\mathrm{id}$ partout ; utilise que $f'-\mathrm{id}$ est encore une dérivée.
\par\smallskip\hfill\small\texttt{Darboux.deriv\_lt\_id\_or\_id\_lt\_deriv}
\end{forme}

\begin{forme}[graphe fermé]\label{f:cg}
$f\in\Darb(\R)$ et graphe de $f$ fermé $\Rightarrow$ $f$ continue.
\par\smallskip\hfill\small\texttt{DarbouxOn.continuous\_of\_isClosed\_graph}
\end{forme}

\begin{forme}[obstruction]\label{f:sgn}
La fonction signe n'est la dérivée d'aucune fonction.
\par\smallskip\hfill\small\texttt{Darboux.not\_isDeriv\_sign}
\end{forme}

\begin{forme}[forme normale, dim.\ $1$]\label{f:norm}
$\rho$ continue $>0$ $\Rightarrow$ $\exists\Phi$ strictement croissante, $\Phi'=\rho$,
$\mathrm{im}\,\Phi$ intervalle, $\int_a^b\rho=\Phi(b)-\Phi(a)$.
\par\smallskip\hfill\small\texttt{Darboux.exists\_chart\_of\_density}
\end{forme}

\section{Trous restants (conjectures candidates)}

Sortie de \texttt{proofspace.py darboux.space} : $42$ paires ouvertes. Les plus parlantes,
avec leur statut :

\begin{center}
\begin{tabular}{@{}lll@{}}
\toprule
Trou & Statut & Commentaire\\
\midrule
$\Darb \Rightarrow \Cont$ & faux & $x^2\sin(1/x)$\\
$\Darb$ stable par $+$ & faux & non formalisé ici\\
$\Darb$ stable par $\circ$ & faux & non formalisé ici\\
$\Der \Rightarrow$ classe de Baire $1$ & vrai (classique) & non formalisé ici\\
$\Darb\ \&\ \text{graphe fermé} \Rightarrow \Cont$ & prouvé & Forme~\ref{f:cg}\\
$\Darb\ \&\ \text{localement borné} \Rightarrow{}?$ & ouvert dans cette carte & à tester\\
\bottomrule
\end{tabular}
\end{center}

\noindent Boucle de travail : (i) fixer un vocabulaire ; (ii) prouver les arêtes évidentes ;
(iii) demander au script les paires et triplets non couverts ; (iv) pour chacun, tenter
l'énoncé \emph{ou} sa négation ; (v) réinjecter. Les Formes~\ref{f:count}, \ref{f:inj},
\ref{f:mono}, \ref{f:fix} sortent de là : ce sont des hyper-arêtes que l'on n'écrit pas
d'habitude, parce que ces théorèmes sont énoncés avec \og{}continue\fg{} à la place de
\og{}Darboux\fg{}.

\section*{Fichiers}

\begin{center}
\begin{tabular}{@{}ll@{}}
\toprule
\texttt{RequestProject/Darboux/Basic.lean} & $\Darb$, ses deux sources, stabilités\\
\texttt{RequestProject/Darboux/Shapes.lean} & Formes \ref{f:sign}--\ref{f:mono}, \ref{f:cg}\\
\texttt{RequestProject/Darboux/Deriv.lean} & Formes \ref{f:cvx}, \ref{f:fix}, \ref{f:sgn}\\
\texttt{RequestProject/Darboux/Space.lean} & le graphe, vérifié\\
\texttt{RequestProject/Darboux/Normal.lean} & Forme \ref{f:norm}\\
\texttt{darboux.space} & entrée de \texttt{proofspace.py}\\
\bottomrule
\end{tabular}
\end{center}

\end{document}
